\section{Weekly Summary}

\begin{tcolorbox}[colback=yellow!10,colframe=black!50,title=AI-Generated Content]
\noindent The summary below was written by Claude (Anthropic) from that week's regression outputs and series descriptions. It has not been reviewed or fact-checked by a human, and should be read as an automated, informal recap -- not analysis.
\end{tcolorbox}

Welcome back to the weekly ritual where we bravely yank two unrelated economic series out of a hat and demand they explain themselves. This week's contestants ranged from portfolio investment assets to Idaho minimum wages to the price of getting places, and, as always, no actual hypotheses were harmed in the making of these regressions.

We kicked things off on November 16 by pairing US Acquisition of Portfolio Investment Assets with a gloriously specific, long-since-discontinued series about Freddie Mac mortgage pools for one-to-four-family homes. The result? A majestic, resounding nothing. The raw regression limped in with an R-squared of 0.013 and a p-value of 0.31, which in plain terms means these two series have never met and would not recognize each other at a party. Detrending somehow made it even worse, dragging the R-squared down to a heroic 0.0007. I have seen stronger relationships between my coffee and my life choices. Honestly, refreshing to open the week with a clean, honest failure instead of a trend mirage.

November 17 is where things got spicy. Idaho's state minimum wage versus Exports of Other Private Services (also discontinued, because apparently this week is a graveyard) produced a raw R-squared of 0.87 -- the kind of number that makes you want to call a press conference. But hold the champagne: the entire point of this project is to not fall for exactly that. And here's the twist -- after detrending, the relationship stubbornly survived, keeping an R-squared of 0.28 and a p-value of 0.0002. That's genuinely rarer than usual around here. Something about Idaho's minimum wage wiggles still correlates with service export wiggles once you strip the shared march of time out of both. Do I have a theory? Sure: both are vaguely tied to a growing economy doing economy things. Is that a real explanation? Absolutely not, and I'd like to formally remind everyone that these are randomly paired series and I am legally and morally obligated to tell you this means nothing.

Finally, November 18 served up the Transportation Services CPI against Retail Sales at Food Services and Drinking Places, with a raw R-squared of 0.95 that looks downright triumphant. As expected, detrending took the air out of the balloon, cutting it to 0.47 -- still significant, but clearly leaning hard on that shared upward drift that everything priced in dollars eventually shares. In other words, prices go up and people spend more at restaurants, both over time, shocking no one who has ever purchased a sandwich. A partial trend-mirage, which is maybe the most on-brand outcome this project can produce. So that's the week: one dud, one suspiciously durable Idaho mystery, and one correlation that's half real and half just the passage of time quietly inflating everything. As always, draw zero conclusions, enjoy the chaos responsibly, and I'll see you next week when the hat spits out more nonsense.
